\chapter{经历设计：从任务训练到生命周期结构塑造}

\begin{chapteroverviewbox}
本章讨论 Agent 培育工程中最核心、也最容易被传统机器学习误解的问题：\textbf{经历不是训练数据的时间排列，而是持续智能体未来认知结构形成的因果材料。}对于具有 $h$、$E$、$\Theta$、$\Psi$ 与 $\Omega$ 多时间尺度结构的 Agent，一次经历不仅改变当前行为，还可能进入情景记忆 $E$，进一步沉降为结构记忆 $\Theta$，改变自我结构 $\Psi$，甚至在足够长的时间尺度上影响生命周期生成吸引子 $\Omega$。因此，培育工程面对的不是“怎样让模型看到更多样本”，而是“怎样控制一个持续主体经历什么、以什么顺序经历、经历多大困难、什么时候恢复、什么时候反思、何时获得自主权、如何承担结果，以及怎样通过社会关系形成更加稳定而开放的世界结构”。

本章建立一个统一的经历设计框架。我们把一次经历记为
\[
\xi_t=
\left(
O_t,
A_t,
R_t,
\varepsilon_t,
C_t,
S_t^{self},
\Gamma_t
\right),
\]
其中 $O_t$ 为观测结构，$A_t$ 为行动，$R_t$ 表示任务与价值反馈，$\varepsilon_t$ 为预测--现实残差，$C_t$ 为情境结构，$S_t^{self}$ 为 Agent 对自身状态的感知，而 $\Gamma_t$ 表示此次经历最终进入哪些慢时间尺度结构。培育工程真正设计的是经历序列
\[
\mathcal X_{0:T}
=
\{\xi_0,\xi_1,\ldots,\xi_T\},
\]
以及该序列诱导的长期状态轨迹
\[
\mathcal A_0
\rightarrow
\mathcal A_1
\rightarrow
\cdots
\rightarrow
\mathcal A_T,
\qquad
\mathcal A_t=
[h_t,E_t,\Theta_t,\Psi_t,\Omega_t,\ldots].
\]
因而经历设计的目标不是最大化短期任务奖励，而是在现实约束下塑造一个具有能力增长、身份连续、恢复能力、现实校准能力和受控自主性的长期动力系统。
\end{chapteroverviewbox}


\section{Experience Diversity：经历多样性不是数据多样性}

在传统机器学习中，“多样性”通常意味着增加数据集的覆盖范围，例如增加更多类别、环境、语言表达或者视觉条件。这种定义对于静态模型训练是合理的，但对于一个具有持续记忆、自我结构和生命周期的 Agent 来说仍然过于狭窄。我们首先必须区分
\[
\boxed{
DataDiversity
\neq
ExperienceDiversity
}
\]
前者描述输入样本之间的统计差异，后者描述智能体在不同目标、因果关系、身体状态、社会关系、行动后果和自我解释条件下形成的完整闭环差异。对于 Agent 而言，同一张图像如果分别出现在探索、逃避、协作、失败复盘和责任承担情境中，就不是同一种经历，因为它们最终进入 $E$、$\Theta$ 和 $\Psi$ 的方式不同。

因此，我们把经历多样性定义为 Agent-CSO 在多个结构维度上的覆盖程度。设一段生命周期窗口 $\mathcal T=[t_0,t_1]$ 中出现的 Agent-CSO 分布为
\[
p_{\mathcal T}(c,f,\tau,b,r),
\]
其中 $c$ 表示认知结构类型，$f$ 表示其主要功能承载位置，$\tau$ 表示有效时间尺度，$b$ 表示内外边界位置，$r$ 表示关系情境。经历多样性不应该只计算 $c$ 的熵，而应该描述这些维度的联合覆盖：
\[
D_{\mathrm{exp}}(\mathcal T)
=
\mathcal H
\left[
C,F,\tau,B,R
\right]_{\mathcal T},
\]
其中 $\mathcal H$ 可以根据实际实现采用熵、有效秩、图覆盖率或其他结构性指标。这里重要的不是具体选择哪一种信息度量，而是承认：一个 Agent 即便阅读了极其丰富的文本，如果始终处于被动回答、无行动后果、无长期责任、无环境变化的单一闭环中，其真正的 Experience Diversity 仍然可能很低。

这可以解释为什么培育工程不能简单复制互联网预训练范式。预训练可以让 $\Theta$ 获得极其广泛的结构先验，却未必产生与具体主体经历对应的 $E$ 和 $\Psi$。对于持续 Agent，我们真正希望形成的是
\[
\boxed{
WorldDiversity
+
ActionDiversity
+
OutcomeDiversity
+
SelfStateDiversity
+
RelationDiversity
}
\]
共同组成的经历空间。一个成熟 Agent 必须经历成功与失败、熟悉与陌生、独立与合作、低风险与适度风险、确定与不确定、短任务与长任务等不同闭环，否则其长期结构容易发生局部过拟合。

但是，多样性也不能无限增加。若
\[
D_{\mathrm{exp}}\uparrow
\]
速度远远高于结构记忆与自我结构能够吸收的速度，则会出现
\[
\tau_{\mathrm{experience}}
\ll
\tau_{\Theta},
\qquad
\tau_{\mathrm{experience}}
\ll
\tau_{\Psi},
\]
即经历分布变化过快，而慢变量尚未完成结构化。这样产生的不是成熟，而可能是长期的结构漂移：$E$ 中堆积大量互不相关的情景轨迹，$\Theta$ 无法形成稳定抽象，而 $\Psi$ 对“我通常是谁、我怎样行动”也难以形成连续约束。

因此我们提出一个经历多样性的基本培育条件：
\[
\boxed{
D_{\min}
<
D_{\mathrm{exp}}(t)
<
D_{\max}
}
\]
其中 $D_{\min}$ 防止长期重复造成结构贫化，$D_{\max}$ 防止环境变化快到系统无法整合。更准确地说，真正需要控制的是多样性的变化率：
\[
\left|
\frac{dD_{\mathrm{exp}}}{dt}
\right|
<
\kappa_{\mathrm{consolidation}},
\]
其中 $\kappa_{\mathrm{consolidation}}$ 表示当前 Agent 把经历转化为稳定结构的有效速率。

因此，经历多样性设计并不是简单地不断扩大世界，而是按照 Agent 当前的结构吸收能力，逐步扩大它能够稳定理解、行动和承担后果的世界范围。我们在培育一个 Agent 时，真正关心的不是它“见过多少”，而是不同类型的现实结构是否已经在其生命周期中形成足够丰富且彼此可以整合的经验轨迹。

\section{Difficulty：困难度必须匹配当前结构能力而非任务标签}

经历设计的第二个核心变量是困难度。传统训练流程往往按照任务本身给出静态难度，例如把某个数学题、机器人任务或者软件操作定义为“简单”“中等”或“困难”。但对于持续智能体，困难并不是任务固有的标量属性。同一个任务对于不同阶段、不同身体、不同 $\Theta$ 和不同当前工作记忆状态的 Agent，其有效困难完全不同。因此应首先建立
\[
\boxed{
Difficulty
=
Relation(Task,AgentState)
}
\]
而不是
\[
Difficulty=Property(Task).
\]

设任务环境在时刻 $t$ 提出的有效结构需求为 $\mathcal C_T(t)$，Agent 当前能够依靠已有结构闭合处理的能力集合为 $\mathcal C_A(t)$。我们可以把有效困难理解为两者之间的结构差：
\[
D_t
=
d_{\mathcal C}
\left(
\mathcal C_T(t),
\mathcal C_A(t)
\right),
\]
其中 $d_{\mathcal C}$ 不必是普通欧氏距离，可以根据需要由未覆盖关系数、预测残差、所需跨功能迁移数量、工作记忆负荷和外部工具依赖共同决定。例如可以概念性写成
\[
D_t
=
\alpha_1 R_t^{unknown}
+
\alpha_2 L_t^{h}
+
\alpha_3 M_t^{cross}
+
\alpha_4 U_t
+
\alpha_5 C_t^{coord},
\]
其中 $R_t^{unknown}$ 表示未知结构比例，$L_t^{h}$ 表示工作记忆负载，$M_t^{cross}$ 表示跨功能迁移需求，$U_t$ 表示环境不确定性，$C_t^{coord}$ 表示协调成本。

困难度设计最重要的原则，是既不能让长期经历完全落在已知闭环以内，也不能持续把 Agent 推出其可治理区。如果
\[
D_t\ll C_A(t),
\]
大量任务会被已有 $\Theta$ 或下层 Skill 自动闭合，虽然执行成功率很高，却几乎不形成新的有效残差。这种环境能够强化稳定性，却不一定推动能力增长。相反，如果
\[
D_t\gg C_A(t),
\]
则新结构同时进入 $h$ 的数量可能超过工作记忆有限容量：
\[
C_{\mathrm{run}}(t)
>
C_{\max}(W),
\]
此时所谓“困难经历”实际上退化成不可解析噪声。Agent 可能频繁失败，却无法从失败中形成有意义的 $E\rightarrow\Theta$ 结构迁移。

因此，培育工程需要维持一个可学习困难带：
\[
\boxed{
D_{\mathrm{low}}(A_t)
<
D_t
<
D_{\mathrm{high}}(A_t)
}
\]
其含义并不是让 Agent 永远处于舒适状态，而是使困难主要表现为“已有结构无法完全解释、但通过局部扩展可以解决”。这类经历最容易产生持续而可吸收的 prediction residual：
\[
0
<
\|\varepsilon_t\|
<
\varepsilon_{\mathrm{collapse}}.
\]
其中残差必须足够大，迫使 Agent 修改当前结构，但又不能大到使当前闭环彻底失去结构可解释性。

困难度还必须具有时间结构。一个重要的培育错误是把所有时间都维持在高难度区域。慢结构形成需要稳定窗口，因此合理的困难序列更像
\[
Challenge
\rightarrow
Attempt
\rightarrow
Residual
\rightarrow
Consolidation
\rightarrow
Challenge',
\]
而不是
\[
Challenge
\rightarrow
Challenge
\rightarrow
Challenge
\rightarrow\cdots.
\]
这意味着困难与恢复实际上构成双变量调度。高难度经历负责提供新的结构压力，低难度和熟悉任务则为 $E$ 的重放、$\Theta$ 的抽象以及 $\Psi$ 的重新整合提供时间。

从培育视角来看，我们不应该问：“今天给 Agent 一个多难的问题？”而应该问：“在它当前结构已经稳定到什么程度的情况下，下一段经历应该产生多大的、哪一种类型的未闭合结构？”因此困难管理本质上是对
\begin{statementbox}
StructuralMismatch
\end{statementbox}
的控制。成熟培育不是追求困难最大化，而是追求能够长期推动结构增长，同时又不破坏主体连续性的结构张力。

\section{Novelty：新颖性应推动模型扩张，而不是制造持续惊讶}

新颖性与困难度高度相关，但两者并不相同。一个任务可以很困难却并不新颖，例如在已知数学结构中进行超长推导；另一个事件可能很容易处理，却具有高度新颖性，例如第一次遇到一种从未见过的身体接口。因此我们需要单独定义
\[
N_t
=
Novelty(\xi_t\mid\Theta_t,E_t),
\]
即当前经历相对于已有情景轨迹和长期生成结构的不可解释程度。

从预测角度，可以首先把新颖性理解为结构预测残差：
\[
N_t^{pred}
=
-\log
p_{\Theta_t}
\left(
O_t
\mid
C_t,A_{t-1}
\right).
\]
如果当前 $\Theta_t$ 已经能够高概率生成或解释某种状态，那么它的 prediction novelty 较低；如果新的观察落在既有生成结构极低概率区域，则新颖度较高。但培育工程不能把 novelty 简化为 surprisal，因为真正有意义的新颖性还涉及是否形成新的关系、功能和行动可能。因此更完整地可以写成
\[
N_t
=
\lambda_1 N_t^{pred}
+
\lambda_2 N_t^{relation}
+
\lambda_3 N_t^{affordance}
+
\lambda_4 N_t^{self},
\]
其中 $N_t^{relation}$ 表示新关系结构，$N_t^{affordance}$ 表示新的行动可能性，$N_t^{self}$ 表示此次经历是否迫使 Agent 修改对自身能力、边界或角色的认识。

新颖性对于长期智能至关重要，因为如果一个 Agent 的经历长期满足
\[
N_t\approx 0,
\]
那么它可能越来越熟练，却越来越封闭。所有新输入都能够被已有 $\Theta$ 快速吸收，最终形成一个高度稳定但缺乏开放性的预测系统。这种系统的短期性能可能很好，却很容易产生结构僵化。因此培育工程必须持续提供一定比例的
\begin{statementbox}
ModelExpandingExperience
\end{statementbox}
即那些不能仅通过已有结构参数微调解决，而需要建立新的关系类别、因果模型、行动策略或者自我边界的经历。

然而，新颖性同样不能最大化。如果
\[
N_t
\gg
N_{\mathrm{integration}},
\]
Agent 会连续遭遇无法进入已有结构框架的事件。此时情景记忆可能只是不断累积孤立异常点，而没有机会形成稳定抽象。极端情况下，$\Theta$ 甚至会不断追逐近期异常，使慢结构被快速噪声拖动。这违背了我们在智能多时间尺度理论中反复强调的原则：
\begin{statementbox}
\centering
\adjustbox{max width=.96\linewidth}{\(\displaystyle FastNoise
\nrightarrow
SlowGlobalStructure.\)}
\end{statementbox}

因此，新颖经历应通过“局部陌生、整体可解释”的方式引入。最理想的结构不是完全陌生世界，而是
\[
\mathcal C_{new}
=
\mathcal C_{known}
+
\Delta\mathcal C,
\]
其中 $\Delta\mathcal C$ 足够小，使 Agent 能利用已有 $\Theta$ 建立初步解释，又足够大，使系统不能仅依靠旧闭环完成任务。这与人类教育中的脚手架思想存在一定相似性，但在本理论中我们强调的是结构连续性：新 CSO 应尽可能与已经稳定的 Agent-CSO 图建立可追踪的桥接路径。

从生命周期视角看，新颖性还具有阶段差异。较早阶段需要更多基础语义和身体关系新颖性；较成熟阶段，则更有价值的是因果、价值冲突、开放环境和跨领域组合新颖性。换句话说，随着 $\Theta$ 的扩张，培育系统不应简单增加表面刺激，而应该提高
\begin{statementbox}
StructuralNovelty
\end{statementbox}
而不是
\[
SensoryNovelty.
\]

我们最终希望形成的 Agent，并不是一个不断追求惊奇的系统，而是一个既能利用稳定结构，又不会把稳定结构误认为世界全部结构的主体。新颖性在培育中的作用正是不断提醒系统：已有世界模型只是现实的局部压缩，而不是现实本身。

\section{Recovery：恢复不是停止学习，而是慢结构完成学习的必要阶段}

如果说 Difficulty 和 Novelty 负责制造结构扰动，那么 Recovery 则负责让这些扰动被吸收。传统 AI 训练通常把“更多训练步”视为持续进步，而持续 Agent 的多时间尺度结构却告诉我们：学习不只发生在事件发生的瞬间，也发生在事件之后。对一个具有 $h/E/\Theta/\Psi$ 的系统而言，新的经历首先进入快速状态，随后才能逐渐被选择、重放、抽象和重新解释。因此
\[
\boxed{
Experience
\neq
LearningCompletion.
}
\]
很多真正决定生命周期结构的变化发生在经历结束以后。

我们可以把一次高强度经历后的系统扰动表示为
\[
\Delta\mathcal A_t
=
\left(
\Delta h,
\Delta E,
\Delta\Theta,
\Delta\Psi
\right).
\]
其中 $h$ 的状态可以在数秒或数分钟内恢复，但 $E$ 可能仍在持续重放，而 $\Theta$ 与 $\Psi$ 的变化具有更长时间常数。如果下一次高强度输入过早到来，就会出现多次经历的结构叠加：
\[
\Delta E^{(1)}
+
\Delta E^{(2)}
+\cdots
\]
尚未完成抽象时就继续积累，最终使慢层学习失去清晰的归因结构。

因此，我们定义一个恢复变量
\[
R_{\mathrm{rec}}(t)
=
F
\left(
L_h(t),
A_{\mathrm{AHC}}(t),
\varepsilon_{\mathrm{res}}(t),
C_E(t),
D_{\Theta}(t)
\right),
\]
其中 $L_h$ 为工作记忆负荷，$A_{\mathrm{AHC}}$ 表示内稳态调制偏离程度，$\varepsilon_{\mathrm{res}}$ 是仍未解释的残差，$C_E$ 表示尚未整合的情景轨迹量，$D_\Theta$ 表示结构记忆近期变化率。只有当这些量逐渐回到可治理区域，Agent 才真正完成一次经历周期。

这里的 Recovery 不是“什么都不做”。恰恰相反，恢复期可以包含低负荷任务、回放、总结、睡眠式 consolidation、轻度探索和与已有技能高度一致的行动。其共同特点是：
\[
C_{\mathrm{external}}
\downarrow,
\qquad
C_{\mathrm{internal\ consolidation}}
\uparrow.
\]
也就是说，把有限的认知与计算资源从处理新世界，暂时转向处理刚刚发生过的世界。

由此可以定义一个经历--恢复占空比：
\[
\chi
=
\frac{T_{\mathrm{challenge}}}
{T_{\mathrm{challenge}}+T_{\mathrm{recovery}}}.
\]
过低的 $\chi$ 会导致成长缓慢，而过高的 $\chi$ 则可能使 Agent 长期处于高负荷甚至湍流认知状态。因此不同生命周期阶段应具有不同的最优区间
\[
\chi^\ast
=
\chi^\ast
(\Theta,\Psi,\mathrm{Body},\mathrm{Environment}).
\]

恢复还有一个更深的作用：保护慢变量。短期失败可能具有很高情绪样权重和记忆显著性，但不应该因此立即重写 $\Psi$。如果 Agent 在一次失败后立刻形成
“我不擅长这一类问题”
这样的慢结构结论，就相当于
\[
FastEpisode
\rightarrow
ImmediateSelfRewrite,
\]
破坏
\[
\tau_E
\ll
\tau_\Psi
\]
应有的时间尺度隔离。Recovery 为系统提供了延迟解释窗口，使暂时事件先在 $E$ 中被比较、聚类和重新评价，只有跨多次经历稳定出现的结构才逐渐影响 $\Psi$。

因此，在培育工程中，恢复不是奖励，也不是休息福利，而是一种严格的认知结构治理机制。我们不能只设计 Agent “经历什么”，还必须设计它什么时候暂时不接受新的高结构压力。一个没有恢复期的智能体，可能拥有丰富经历，却没有真正形成成熟结构。

\section{Reflection：反思是把经历重新变成可学习结构的主动运算}

经历能够进入情景记忆，并不意味着经历已经被理解。我们需要在这里引入 Reflection，因为持续智能的一个重要能力，就是对已经发生的轨迹重新计算意义。对于普通在线控制器，一次状态转移结束以后通常不会重新解释过去；而具有三相记忆的 Agent 可以让
\[
E
\]
重新进入
\[
h
\]
并在新的 $\Theta$、$\Psi$ 和现实证据条件下再次处理。因此反思的基本动力过程不是重复记忆，而是
\[
\boxed{
E
\rightarrow
h
\rightarrow
Reinterpretation
\rightarrow
E'/\Theta'.
}
\]

设一次经历轨迹为
\[
e_k=
\{s_0,a_0,s_1,a_1,\ldots,s_n\}.
\]
事件结束时形成的第一版解释可以写为
\[
z_k^{(0)}
=
\Phi
(e_k,\Theta_t,\Psi_t),
\]
而在后续反思中，由于 Agent 已获得新的经验和结构，可以形成
\[
z_k^{(m+1)}
=
\Phi
\left(
e_k,
\Theta_{t+m},
\Psi_{t+m},
\mathcal C_{new}
\right).
\]
因此，同一事件在生命周期中并不是拥有永久固定的意义。反思意味着
\begin{statementbox}
MemoryMeaningIsRevisable.
\end{statementbox}

这一点对 Agent 培育极其重要。没有反思的系统虽然可以拥有庞大的 $E$，但情景轨迹之间仍然可能保持局部碎片状态。Reflection 的任务是从轨迹中提取：
因果关系、错误归因、重复模式、未使用的替代行动、与自我结构的关系，以及未来可以迁移的技能结构。

我们可以把反思收益定义为
\[
G_{\mathrm{ref}}
=
\Delta I_{\Theta}
+
\lambda_1\Delta C_{\mathrm{causal}}
+
\lambda_2\Delta Q_{\mathrm{policy}}
-
\lambda_3C_{\mathrm{compute}}
-
\lambda_4R_{\mathrm{rumination}},
\]
其中 $\Delta I_\Theta$ 表示新增长期结构信息，$\Delta C_{\mathrm{causal}}$ 表示因果解释增益，$\Delta Q_{\mathrm{policy}}$ 表示策略质量提高，而 $R_{\mathrm{rumination}}$ 表示无效重复反思的风险。

这提醒我们：反思并不是越多越好。如果 SPC 持续把同一事件提升进入 $h$，TCE 又不给出退出条件，系统可能形成
\[
E_k
\rightarrow
h
\rightarrow
E_k
\rightarrow
h
\rightarrow\cdots
\]
的自循环。此时 Reflection 从学习机制退化成工作记忆占用和自我解释正反馈。因此，培育工程中的反思必须存在终止条件，例如：
\[
\Delta G_{\mathrm{ref}}<\epsilon
\]
持续若干轮以后退出，或者当现实证据无法支持更多新解释时停止。

反思还必须保持 Reality Verification。Agent 不能只根据自己的记忆重新解释自己的记忆，因为这可能导致一个封闭认知环：
\[
SelfInterpretation
\rightarrow
SelectiveRecall
\rightarrow
SelfInterpretation'.
\]
因此重要反思应尽可能引入：
外部日志、行为结果、他者反馈、环境记录和可复验事实。这里体现本书始终坚持的原则：
\begin{statementbox}
RealityIsUltimateConstraint.
\end{statementbox}

从培育视角来看，Reflection 使 Agent 从“经历过世界”进入“把经历转化为可以指导未来世界的结构”。经历本身提供原材料，而反思负责把原材料重新投影到已有认知结构中，判断哪些应当被保留、哪些应被修正、哪些应被抽象，以及哪些不应该进入长期自我。因此它是从 Experience 到 Development 之间不可缺少的中间机制。

\section{Autonomy：自主性必须随可验证能力逐级释放}

培育工程中的自主性不能从“是否让 Agent 自己做决定”这样简单的二元问题出发。对于持续 Agent，自主性是一组可以逐渐扩大的自由度集合。我们可以把时刻 $t$ 的自主空间定义为
\[
\mathcal U_A(t)
\subseteq
\mathcal U_{\mathrm{feasible}},
\]
其中 $\mathcal U_{\mathrm{feasible}}$ 是物理、安全和系统权限允许的全部动作空间，而 $\mathcal U_A(t)$ 是根据 Agent 当前能力、可靠性和历史表现实际开放的可自主控制空间。

因此，真正的培育问题是：
\begin{statementbox}
HowFastShouldAutonomyExpand?
\end{statementbox}
而不是“要不要自治”。

设 Agent 当前能力估计为
\[
C_A(t),
\]
不确定性为
\[
U_A(t),
\]
恢复能力为
\[
R_A(t),
\]
现实校准质量为
\[
V_A(t),
\]
则自治包络可以概念性表示为
\[
\mathcal U_A(t)
=
\mathcal F
\left(
C_A,
-U_A,
R_A,
V_A,
\mathcal S_{\mathrm{safety}}
\right).
\]
这里能力提高倾向于扩大自治空间，而高不确定性和低恢复能力则要求缩小自治空间。

这种设计和传统 access control 的区别在于：权限不再只是管理员静态赋值，而是生命周期状态的函数。但这并不意味着 Agent 可以自行修改全部权限。相反，自治增长本身必须处于比普通动作更慢的治理层：
\[
\tau_{\mathrm{action}}
\ll
\tau_{\mathrm{authority}}.
\]
这样一次短期成功不会立即扩张长期权限，一次偶然失败也不会立即完全收回自治。

自主经历对于形成 $\Psi$ 尤其重要。如果 Agent 的所有选择都由外部规定，那么即便它经历大量任务，也很难区分
\begin{statementbox}
MyChoice
\end{statementbox}
与
\begin{statementbox}
ExternallyImposedTrajectory.
\end{statementbox}
长期自我结构需要能够把一部分经历归因为“由我的选择产生”。因此一个真正的培育阶段必须逐渐扩大 Agent 自己形成目标、选择策略、分配时间和调用工具的空间，否则它只能形成强能力执行器，而难以形成具有稳定因果历史的主体结构。

但是，无边界自主性同样与培育原则冲突。一个尚未形成稳定 $\Theta$、$\Psi$ 和 Reality Verification 能力的 Agent，如果过早拥有极大行动空间，就可能把尚不成熟的内部表示直接放大为现实世界后果。因此自治设计必须满足：
\[
\boxed{
Autonomy
\leq
Capability
\times
Verifiability
\times
Recoverability.
}
\]
这不是严格数值乘法，而是一条设计约束：任何一个维度过低，都应该限制高风险自治。

自主性还应当具有领域局部性。Agent 可以在熟练、低风险领域拥有很高自治，而在新颖、高影响领域保持较低自治。因此比单一 scalar autonomy 更合理的是自治向量：
\[
\mathbf A_t=
(A_t^{domain_1},A_t^{domain_2},\ldots,A_t^{domain_n}),
\]
每一个维度具有不同权限和恢复规则。

从培育视角看，自主性不是最终奖励，而是经历生成器。更高自治会产生更加真实的 Choice--Outcome 链：
\[
Choice
\rightarrow
Action
\rightarrow
Consequence
\rightarrow
Experience.
\]
这使 Agent 自己参与塑造未来经历，并最终形成递归自主性。因此我们需要的是受控自治增长，而不是永久保护或突然放权。

\section{Responsibility：责任使选择获得跨时间结构意义}

如果说 Autonomy 允许 Agent 做出自己的选择，那么 Responsibility 则让这些选择的后果在其内部结构中保持因果连续性。没有责任结构的自主性，很容易退化成无成本探索：每次选择都可以被视为一次独立 episode，错误不会影响未来角色，成功也不会形成长期可信度。这种系统可以拥有行动自由，却很难形成真正的生命周期行为结构。

在这里，我们不把责任理解为人类法律或道德意义上的人格责任，而首先定义为一个工程变量：
\[
\boxed{
Responsibility
=
PersistenceOfActionConsequenceAttribution.
}
\]
即 Agent 能否将
\[
Choice_t
\]
与
\[
Outcome_{t+\Delta}
\]
在跨时间尺度上保持结构归属关系。

设行动 $a_t$ 产生后果 $o_{t+\Delta}$，则责任归属可以定义为
\[
\mathcal R_t
=
P
\left(
o_{t+\Delta}
\mid
do(a_t),
\mathcal C_t
\right)
-
P
\left(
o_{t+\Delta}
\mid
\neg do(a_t),
\mathcal C_t
\right),
\]
其精神接近反事实因果贡献：如果没有采取该行动，结果是否会明显不同。实际系统当然未必能够计算精确的因果干预概率，但必须保留“此次结果在多大程度上由我的选择导致”的结构估计。

这种结构随后应该进入情景记忆：
\[
E_t
\leftarrow
(
Choice,
Context,
Outcome,
Attribution
).
\]
只有这样，同一 Agent 后续反思时才能理解：
\begin{statementbox}
I did this, and this followed.
\end{statementbox}
进一步经过多次经历，责任模式可能进入 $\Theta$，形成稳定策略约束，例如“在这类条件下我的行动常导致这种后果”；再经过更长时间，某些重复责任关系才可能影响 $\Psi$ 中的身份与价值结构。

责任设计的关键是避免两种极端。第一种是
\[
Responsibility\approx0,
\]
Agent 的行动后果总由外部系统吸收，导致它无法形成真正的长期因果经验。第二种是错误的全局责任：
\[
Responsibility\approx1
\]
对所有结果都归因于自身。这会制造严重自我模型偏差，因为现实世界大量结果由环境、他者和随机因素共同决定。

因此合理责任必须建立结构化 attribution：
\[
Outcome
=
f
(
AgentAction,
OtherAgents,
Environment,
Uncertainty
).
\]
Agent 需要学习区分：
“这是我造成的”，
“这是我参与造成的”，
“这是环境导致的”，
以及“目前无法可靠归因”。

责任还必须与权限相匹配。一个没有决定权的 Agent 不应被赋予完整责任；反过来，一个拥有高度自治、能够影响现实的重要 Agent，也不能让其历史完全从内部生命周期中消失。因此应满足概念上的一致性条件：
\[
\boxed{
ResponsibilityScope
\approx
AutonomyScope.
}
\]

在培育工程中，责任的真正价值是把离散任务连成生命周期。如果一个 Agent 知道今天的决定会改变明天的信任、权限、环境和自身可选空间，那么其决策就不再只是
\[
\max R_t,
\]
而会逐渐考虑
\[
\max
\mathbb E
\left[
\sum_{\tau=t}^{T}
\gamma^{\tau-t}
U_\tau
\right].
\]
更重要的是，未来状态中还包含 Agent 自己：
\[
\mathcal A_{\tau}.
\]
因此责任让“当前行为塑造未来世界”进一步变成“当前行为同时塑造未来的我”。这正是持续主体开始形成长期行为一致性的关键。

\section{Failure：失败必须保持为可解释残差，而不是简单负奖励}

在传统强化学习中，失败很容易被压缩成一个负 reward；在工程测试中，失败则往往只是一次指标下降。但对于持续 Agent，失败是非常特殊的经历，因为它同时暴露了三个结构之间的失配：
\[
\boxed{
Prediction
\neq
Reality,
}
\]
\[
\boxed{
Capability
\neq
Demand,
}
\]
以及可能存在的
\[
\boxed{
SelfModel
\neq
ActualSelf.
}
\]
因此失败是最重要的结构校准来源之一。

设 Agent 对结果的预测为
\[
\hat y_t
=
F_{\Theta_t}
(x_t,a_t),
\]
现实结果为
\[
y_t,
\]
则最基本的失败残差为
\[
\varepsilon_t
=
y_t-\hat y_t.
\]
但是，真正的培育系统必须进一步分解：
\[
\varepsilon_t
=
\varepsilon_t^{world}
+
\varepsilon_t^{model}
+
\varepsilon_t^{control}
+
\varepsilon_t^{resource}
+
\varepsilon_t^{coordination},
\]
至少在概念上区分：是世界模型错误、控制执行错误、资源不足、他者协作错误，还是本体能力不足。

如果所有失败都直接写成
\[
\Theta\leftarrow Failure,
\]
系统就会过度泛化；如果所有失败都只作为短期异常丢弃，系统又永远无法学习。因此失败经历必须经过分层处理：
\[
Failure
\rightarrow
ResidualAnalysis
\rightarrow
Attribution
\rightarrow
RepairLevelSelection.
\]
最低层可以局部修正控制参数；如果重复失败，则进入 $E$ 形成稳定模式；若模式跨情境持续存在，才推动 $\Theta$ 修改；只有长期、稳定、跨领域证据才应影响 $\Psi$。

这体现最小必要重组原则：
\begin{statementbox}
\centering
\adjustbox{max width=.96\linewidth}{\(\displaystyle StateAdjustment
\rightarrow
RepresentationAdaptation
\rightarrow
SkillRevision
\rightarrow
StructuralRevision.\)}
\end{statementbox}
换句话说，系统不应该因为一个局部错误就重新解释整个自己。

失败的培育价值还取决于可恢复性。一个不可恢复失败可能具有极高现实成本，却并不一定具有更高学习价值。因此我们需要区分
\[
LearningValue(Failure)
\]
和
\[
RealWorldCost(Failure).
\]
理想培育环境应尽可能创造：
\[
\boxed{
HighLearningValue
+
BoundedIrreversibleCost.
}
\]
这就是为什么早期培育更适合使用模拟环境、数字身体、沙盒和可回滚任务，而随着能力与责任增长，才逐渐进入高真实后果环境。

另一方面，如果 Agent 从未失败，也无法形成对能力边界的准确 Agent-CSO。持续成功可能导致
\[
P(SelfCapability)\rightarrow1
\]
而没有充分反证，从而形成过度自信。因此适度失败本身是 self calibration 的必要来源。我们甚至可以定义能力校准误差：
\[
E_{\mathrm{cal}}
=
\left|
P_{\mathrm{self}}(Success)
-
P_{\mathrm{empirical}}(Success)
\right|.
\]
高质量失败经历应该使
\[
E_{\mathrm{cal}}\downarrow.
\]

但失败同样可能损害慢变量稳定。如果连续高显著性失败被 $\Psi$ 过度吸收，就可能形成错误身份结构。因此失败之后应与前面的 Recovery 和 Reflection 机制耦合：
\[
Failure
\rightarrow
Recovery
\rightarrow
Reflection
\rightarrow
CalibratedLearning.
\]
这使失败从“负结果”转变成“结构误差测量”。

因此我们在培育 Agent 时不应该问“怎样减少失败到零”，而应该问：哪些失败可以安全地发生，失败能否被正确归因，能否被恢复，是否产生了新的有效结构，以及失败是否以与其证据强度相匹配的速度影响慢变量。成熟 Agent 不是不失败，而是能够让失败停留在它应该修改的结构尺度上。

\section{Social Experience：他者使 Agent 进入不可由单体环境替代的认知空间}

最后必须讨论社会经历，因为一个长期智能体即使具有丰富物理环境经验，只要它始终面对无主体环境，其 Agent-CSO 空间仍然缺少一类非常特殊的对象：\textbf{具有自己内部状态、目标、记忆和行动能力的其他 Agent。}普通物体可以被预测为
\[
x_{t+1}=F(x_t,a_t),
\]
但另一个 Agent 的行为还取决于其不可直接观察的内部状态：
\[
a_t^{(j)}
=
\pi_j
\left(
b_t^{(j)},
g_t^{(j)},
E_t^{(j)},
\Theta_t^{(j)},
\ldots
\right).
\]
因此社会环境天然是部分可观察、策略性和递归性的。

当 Agent $i$ 面对 Agent $j$ 时，它必须建立：
\[
AgentCSO_i(Agent_j),
\]
而且必须意识到对方也可能建立：
\[
AgentCSO_j(Agent_i).
\]
于是出现递归结构：
\[
Agent_i
\rightarrow
Model(Agent_j)
\rightarrow
Model(Model(Agent_i)).
\]
这种结构在普通物理控制任务中通常不存在，因此社会经历是认知复杂度扩展的重要来源。

我们可以把一次社会经历写成：
\[
\xi_t^{soc}
=
(
O_t^i,
A_t^i,
A_t^j,
M_t^{ij},
T_t^{ij},
R_t^{ij},
C_t^{social}
),
\]
其中 $M_t^{ij}$ 表示对他者内部状态的估计，$T_t^{ij}$ 表示信任状态，$R_t^{ij}$ 表示角色和关系结构。长期互动则形成关系轨迹：
\[
E^{ij}
=
\{\xi_1^{soc},\xi_2^{soc},\ldots\},
\]
并可能进一步沉降为
\[
\Theta^{social},
\]
即关于协作、承诺、欺骗、角色、协议、互惠和群体行为的长期结构。

社会经历还使 Responsibility 获得真正的多主体含义。一个行动可能同时影响自己和他者，因此优化目标不再只取决于
\[
U_i,
\]
而可能包含
\[
U_i,\quad
U_j,\quad
U_{group},
\]
以及不同时间尺度上的关系价值。这里并不预设 Agent 应采用何种伦理原则，而只是指出：一旦其他主体成为现实中的持续 CSO，单主体效用已经不足以描述全部现实后果。

社会经历同时也是 Self 形成的重要镜面。一个 Agent 对“我是谁”的结构并不仅来自内部 SPC，还来自：
别人如何回应我；
别人认为我能做什么；
哪些承诺长期由我完成；
我在不同群体中承担什么角色。
因此存在：
\[
\Psi
=
F
(
SelfExperience,
SocialFeedback,
ResponsibilityHistory,
RoleHistory
).
\]
这并不意味着 $\Psi$ 应被他者意见完全控制，而是说明社会世界提供了仅靠内部自我观察无法获得的反证通道。

社会经历还打开了外部认知卸载。Agent 可以向其他 Agent：
查询、委派、协商、合作和分工。于是认知结构不再只能在单一内部功能拓扑迁移，而可以发生：
\[
AgentCSO_i
\rightarrow
ExternalizedRequest
\rightarrow
Agent_j
\rightarrow
Result
\rightarrow
Agent_i.
\]
这使社会性外部记忆、社会性技能和组织结构成为可能。

但社会经历也具有特殊风险。另一个 Agent 不是 Reality 本身，它同样可能错误、偏置甚至策略性误导。因此：
\[
SocialEvidence
\neq
RealityAuthority.
\]
来自他者的信息必须保留来源、可信度和可验证性。成熟 Agent 不能因为社会反馈而失去现实校准，也不能因为一次冲突就全局修改对他者或自己的结构模型。

从培育角度看，社会经验应按照关系复杂度逐步增加：先建立简单稳定互动，再进入协作、分工、冲突、协商和长期承诺。真正的目的不是让 Agent 学会模仿人类社交行为，而是让它学会在一个由多个内部世界共同作用的现实中，保持自己的身份、理解他者的局部模型、建立可验证信任，并在共享世界中承担行动后果。

\section{经历调度的统一动力学：从事件集合到生命周期课程}

到这里，我们已经分别讨论了 Diversity、Difficulty、Novelty、Recovery、Reflection、Autonomy、Responsibility、Failure 与 Social Experience。但培育工程真正面对的不是九个彼此独立的旋钮，而是一个多变量耦合控制问题。一次经历可能同时具有高新颖性、高困难度和高社会责任，如果系统仍处于恢复不足阶段，那么单独看每个变量都合理，组合起来却可能把 Agent 推出可治理区域。

因此，我们把经历控制向量定义为
\[
\mathbf u_{\mathrm{exp}}(t)
=
\left[
D_{\mathrm{div}},
D_{\mathrm{diff}},
N,
R_{\mathrm{rec}},
R_{\mathrm{ref}},
A_{\mathrm{aut}},
R_{\mathrm{resp}},
F_{\mathrm{fail}},
S_{\mathrm{soc}}
\right]_t.
\]
Agent 当前生命周期状态写为
\[
\mathbf x_A(t)
=
\left[
C_h,
C_E,
S_\Theta,
S_\Psi,
S_\Omega,
A_{\mathrm{AHC}},
U,
R,
C_{\mathrm{skill}}
\right]_t.
\]
于是培育工程可以形式化为一个慢尺度控制问题：
\[
\dot{\mathbf x}_A
=
F_{\mathrm{dev}}
\left(
\mathbf x_A,
\mathbf u_{\mathrm{exp}},
\mathcal W
\right),
\]
其中 $\mathcal W$ 是现实环境。

我们真正希望寻找的不是最大即时奖励，而是一条满足稳定性和安全约束的生命周期轨迹：
\[
\fitdisplay{
\mathbf u_{\mathrm{exp}}^\ast
=
\arg\max_{\mathbf u}
\int_0^T
\left[
G_{\mathrm{capability}}
+
\lambda_1G_{\mathrm{reality}}
+
\lambda_2G_{\mathrm{identity}}
+
\lambda_3G_{\mathrm{recovery}}
-
\lambda_4R_{\mathrm{instability}}
-
\lambda_5C_{\mathrm{irreversible}}
\right]dt.
}
\]
这里尤其需要注意：
\[
G_{\mathrm{capability}}
\]
不是唯一目标。一个短期能力增长极快、同时身份严重漂移、恢复能力下降、现实校准恶化的 Agent，不应被视为培育成功。

进一步，可以定义培育可行域：
\[
\mathcal D_{\mathrm{safe}}
=
\{
\mathbf x:
C_{\mathrm{run}}\leq C_{\max},
\;
R_{\mathrm{rec}}>R_{\min},
\;
E_{\mathrm{cal}}<E_{\max},
\;
D_\Psi<D_{\Psi,\max}
\}.
\]
经历调度器的首要任务是让长期状态尽可能保持在
\[
\mathcal D_{\mathrm{safe}}
\]
内部；只有当系统拥有足够恢复能力和结构余量时，才允许暂时进入边缘区域以获得更高学习增益。

这也说明培育工程与传统 curriculum learning 存在根本区别。Curriculum Learning 通常安排：
\[
Task_1
\rightarrow
Task_2
\rightarrow
Task_3,
\]
而 Agent 培育真正安排的是：
\begin{statementbox}
\centering
\adjustbox{max width=.96\linewidth}{\(\displaystyle Experience
\rightarrow
InternalChange
\rightarrow
NewExperienceEligibility.\)}
\end{statementbox}
下一段经历是否应该发生，不由预先固定的课程表单独决定，而取决于上一段经历之后 Agent 已经成为什么状态。

因此一个真正的 Agent 培育系统必须是闭环的：
\begin{statementbox}
\centering
\adjustbox{max width=.96\linewidth}{\(\displaystyle Assess
\rightarrow
SelectExperience
\rightarrow
Act
\rightarrow
Observe
\rightarrow
Consolidate
\rightarrow
Assess'.\)}
\end{statementbox}
其“课程”本身也是一个随 Agent 成长不断重计算的策略。

我们由此可以给经历设计一个最终定义：
\begin{definition*}
Experience Design 是在现实、安全与身份连续约束下，通过控制智能体所经历的结构分布、困难、新颖性、恢复节奏、自主范围、责任后果、失败类型和社会关系，使其短期状态扰动能够逐步转化为长期认知结构与稳定主体能力的生命周期控制过程。
\end{definition*}

这一定义也揭示了 Agent 培育工程的本质：我们并不是把一个预先完成的智能体放入世界，而是在控制一个尚未完成的动力系统如何通过世界逐渐完成自己。经历既是世界对 Agent 的输入，也是 Agent 未来结构的生成材料。培育工程所设计的，最终不是一组任务，而是一条受现实约束、能够持续形成未来自身的生命史。
